% ============================================================
%  ANNEXES
% ============================================================
\chapter{Extraits de code significatifs}
\label{annexe:code}

\section{Configuration complète des variables d'environnement}

\begin{lstlisting}[language=bash, caption={backend/.env -- Variables d'environnement complètes}, label={lst:env_complet}]
# Serveur
PORT=5000
NODE_ENV=development

# Base de données
MONGO_URI=mongodb://localhost:27017/medical_anonymizer

# Authentification JWT
JWT_SECRET=votre_cle_secrete_jwt_tres_longue_32_caracteres_minimum
JWT_EXPIRE=7d

# Clé d'enregistrement administrateur
ADMIN_REGISTRATION_KEY=PURA_ADMIN_2026

# URL du pipeline IA FastAPI
FASTAPI_URL=http://localhost:8000

# Frontend (CORS)
FRONTEND_URL=http://localhost:3000

# MinIO Object Storage
MINIO_ENDPOINT=localhost:9000
MINIO_ACCESS_KEY=minioadmin
MINIO_SECRET_KEY=minioadmin
MINIO_BUCKET=anonymized-images
MINIO_SECURE=false
\end{lstlisting}

\section{Modèle Mongoose History (journal des opérations)}

\begin{lstlisting}[language=JavaScript, caption={Modèle Mongoose History}, label={lst:history_model}]
// backend/src/models/History.js
const mongoose = require('mongoose')

const HistorySchema = new mongoose.Schema({
  user: {
    type: mongoose.Schema.Types.ObjectId,
    ref: 'User', required: true
  },
  filename:        { type: String, required: true },
  status:          { type: String, enum: ['success','failed'],
                     default: 'success' },
  category:        String,
  confidence:      Number,
  regionsDetected: Number,
  regionsRedacted: Number,
  tagsAnonymized:  Number,
  processingTime:  Number,
  minioUri:        String,
  downloadUrl:     String,
  confThreshold:   Number,
  padding:         Number,
  borderMargin:    Number,
  borderPct:       Number
}, { timestamps: true })

module.exports = mongoose.model('History', HistorySchema)
\end{lstlisting}

\section{Route d'anonymisation backend}

\begin{lstlisting}[language=JavaScript, caption={Route d'anonymisation Node.js}, label={lst:route_anon}]
// backend/src/routes/anonymize.js
const express    = require('express')
const router     = express.Router()
const { protect } = require('../middleware/auth')
const { upload }  = require('../middleware/upload')
const { anonymizeImage } = require('../controllers/anonymizeController')

router.post('/', protect, upload.single('file'), anonymizeImage)

module.exports = router
\end{lstlisting}

\section{Middleware d'upload Multer}

\begin{lstlisting}[language=JavaScript, caption={Configuration Multer -- stockage mémoire et filtre}, label={lst:multer}]
// backend/src/middleware/upload.js
const multer = require('multer')

const storage = multer.memoryStorage()

const fileFilter = (req, file, cb) => {
  const allowed = ['.jpg', '.jpeg', '.png', '.bmp',
                   '.tiff', '.tif', '.dcm', '.dicom']
  const ext = '.' + file.originalname.split('.').pop().toLowerCase()
  if (allowed.includes(ext)) {
    cb(null, true)
  } else {
    cb(new Error(`Format non supporté: ${ext}`), false)
  }
}

exports.upload = multer({
  storage,
  fileFilter,
  limits: { fileSize: 50 * 1024 * 1024 }  // 50 Mo
})
\end{lstlisting}

\section{Anonymiseur de métadonnées DICOM}

\begin{lstlisting}[language=Python, caption={MetadataAnonymizer -- suppression des tags PHI DICOM}, label={lst:metadata}]
# anonymizer/metadata_anonymizer.py (extrait)
PHI_TAGS = [
    (0x0010, 0x0010),  # Patient Name
    (0x0010, 0x0020),  # Patient ID
    (0x0010, 0x0030),  # Patient Birth Date
    (0x0010, 0x0040),  # Patient Sex
    (0x0010, 0x1010),  # Patient Age
    (0x0010, 0x1030),  # Patient Weight
    (0x0008, 0x0080),  # Institution Name
    (0x0008, 0x0081),  # Institution Address
    (0x0008, 0x0090),  # Referring Physician Name
    (0x0008, 0x1048),  # Physician(s) of Record
    (0x0008, 0x1050),  # Performing Physician Name
    (0x0020, 0x0010),  # Study ID
    (0x0008, 0x0020),  # Study Date
    (0x0008, 0x0030),  # Study Time
    (0x0008, 0x103E),  # Series Description
    (0x0040, 0xA124),  # UID (de-identify)
]

class MetadataAnonymizer:
    def anonymize(self, dataset):
        count = 0
        for tag in PHI_TAGS:
            if tag in dataset:
                try:
                    dataset[tag].value = ''
                    count += 1
                except Exception:
                    pass
        return dataset, count
\end{lstlisting}

\section{Endpoint FastAPI complet /anonymize}

\begin{lstlisting}[language=Python, caption={Endpoint FastAPI /anonymize -- orchestration du pipeline}, label={lst:fastapi_endpoint}]
# api/main.py -- endpoint /anonymize (extrait condensé)
@app.post("/anonymize")
async def anonymize_image(
    file: UploadFile = File(...),
    conf_threshold: float = Form(0.1),
    padding:        int   = Form(5),
    border_margin:  int   = Form(100),
    border_pct:     float = Form(0.20)
):
    start = time.time()
    # Sauvegarde temporaire
    temp_in = TEMP_DIR / f"tmp_{datetime.now():%Y%m%d_%H%M%S}_{file.filename}"
    content = await file.read()
    temp_in.write_bytes(content)

    # Stage 1: Classification
    classifier = MedicalImageClassifier()
    category, confidence, _ = classifier.classify_image(str(temp_in))
    if "non_medical" in category:
        raise HTTPException(400, f"Image non médicale ({confidence:.2f})")

    # Stage 2-3: Validation + Métadonnées DICOM
    validator = ImageValidator()
    result    = validator.validate(str(temp_in))
    is_dicom  = result.is_dicom
    pixels    = result.pixel_array
    dataset   = result.dataset
    tags_anon = 0
    if is_dicom:
        meta = MetadataAnonymizer()
        dataset, tags_anon = meta.anonymize(dataset)

    # Stage 4: CLAHE
    prep     = BorderPreprocessor(border_pct=0.15)
    enhanced = prep.enhance(pixels if pixels is not None else np.array([]))

    # Stage 5: OCR double moteur
    paddle = TextDetector(lang="en", conf_threshold=conf_threshold)
    easy   = EasyTextDetector(conf_threshold=conf_threshold,
                               border_pct=border_pct)
    p_reg  = paddle.detect_text(enhanced)
    e_reg  = easy.detect_text(enhanced)
    merged = merge_and_deduplicate(p_reg, e_reg, 0.5)

    # Stage 6: Redaction
    redactor = PixelRedactor(padding=padding, border_margin=border_margin)
    out_data, n_redacted = redactor.redact(
        dataset if is_dicom else pixels,
        merged, padding=padding,
        border_margin=border_margin, redact_all_regions=True
    )

    # Stage 7: Sauvegarde + MinIO
    out_file = OUTPUT_DIR / f"anonymized_{file.filename}"
    # ... sauvegarde + upload MinIO ...

    return JSONResponse({
        "status": "success",
        "category": category, "confidence": float(confidence),
        "regions_detected": len(merged),
        "regions_redacted": n_redacted,
        "tags_anonymized": tags_anon,
        "processing_time": round(time.time() - start, 2),
        "filename": out_file.name
    })
\end{lstlisting}

% ============================================================
\chapter{Documentation de l'API}
\label{annexe:api}

\section{Swagger -- Documentation automatique FastAPI}

FastAPI génère automatiquement une documentation interactive Swagger UI accessible à l'adresse \texttt{http://localhost:8000/docs}. Cette documentation liste tous les endpoints, les schémas de requête/réponse et permet de tester l'API directement depuis le navigateur.

\begin{figure}[H]
  \centering
  \includegraphics[width=\textwidth]{example-image}
  \caption{Documentation Swagger auto-générée par FastAPI (http://localhost:8000/docs)}
  \label{fig:swagger}
\end{figure}

\section{Exemples de requêtes et réponses}

\begin{lstlisting}[language=Python, caption={Exemple de requête d'anonymisation avec curl}, label={lst:curl}]
curl -X POST http://localhost:5000/api/anonymize \
  -H "Authorization: Bearer <JWT_TOKEN>" \
  -F "file=@/path/to/xray.jpg" \
  -F "conf_threshold=0.1" \
  -F "padding=5" \
  -F "border_margin=100" \
  -F "border_pct=0.20"
\end{lstlisting}

\begin{lstlisting}[language=Python, caption={Exemple de réponse JSON en cas de succès}, label={lst:response_json}]
{
  "success": true,
  "data": {
    "status": "success",
    "filename": "anonymized_xray.jpg",
    "category": "chest",
    "confidence": 0.9412,
    "regions_detected": 8,
    "regions_redacted": 7,
    "regions_skipped": 1,
    "tags_anonymized": 0,
    "processing_time": 2.84,
    "download_url": "http://localhost:9000/anonymized-images/chest/2026/04/08/anonymized_xray.jpg?...",
    "format": "JPEG"
  }
}
\end{lstlisting}

\begin{lstlisting}[language=Python, caption={Exemple de réponse en cas d'image non médicale}, label={lst:response_error}]
{
  "status": "failed",
  "error": "non_medical",
  "message": "Image classified as non-medical (confidence: 0.87)",
  "classification": "non_medical",
  "confidence": 0.87
}
\end{lstlisting}

% ============================================================
\chapter{Guide d'installation et de déploiement}
\label{annexe:install}

\section{Prérequis système}

\begin{itemize}
  \item \textbf{Système d'exploitation}~: Windows 10/11, Ubuntu 20.04+ ou macOS 12+
  \item \textbf{Python}~: Version 3.10.x (recommandé~: 3.10.11)
  \item \textbf{Node.js}~: Version 18.x ou supérieure
  \item \textbf{MongoDB}~: Version 6.0 ou supérieure (local ou Atlas)
  \item \textbf{Docker}~: Version 24.0 ou supérieure (pour MinIO)
  \item \textbf{RAM}~: Minimum 8 Go (16 Go recommandé pour CLIP + PaddleOCR)
  \item \textbf{Espace disque}~: Minimum 10 Go pour les modèles et dépendances
\end{itemize}

\section{Installation étape par étape}

\begin{lstlisting}[language=bash, caption={Procédure d'installation complète}, label={lst:install}]
# 1. Cloner le projet
git clone https://github.com/[votre-repo]/PFE_Test.git
cd PFE_Test

# 2. Créer et activer l'environnement virtuel Python
python -m venv venv
# Windows:
venv\Scripts\activate
# Linux/macOS:
source venv/bin/activate

# 3. Installer les dépendances Python
pip install -r requirements.txt

# 4. Installer les dépendances Node.js (backend)
cd backend && npm install && cd ..

# 5. Installer les dépendances Node.js (frontend)
cd client && npm install && cd ..

# 6. Configurer les variables d'environnement
cp backend/.env.example backend/.env
# Éditer backend/.env avec vos valeurs

# 7. Démarrer MinIO avec Docker
docker compose up -d minio

# 8. Démarrer MongoDB (si local)
mongod --dbpath /data/db

# 9. Démarrer le backend Node.js
cd backend && npm start &

# 10. Démarrer le pipeline FastAPI
python -m uvicorn api.main:app --reload --port 8000 &

# 11. Démarrer le frontend React
cd client && npm run dev
\end{lstlisting}

\section{Accès à l'application}

Une fois tous les services démarrés~:

\begin{itemize}
  \item \textbf{Frontend React}~: \url{http://localhost:3000}
  \item \textbf{Backend Node.js}~: \url{http://localhost:5000}
  \item \textbf{API FastAPI + Swagger}~: \url{http://localhost:8000/docs}
  \item \textbf{Console MinIO}~: \url{http://localhost:9001} (minioadmin / minioadmin)
\end{itemize}

% ============================================================
\chapter{Glossaire des termes techniques}
\label{annexe:glossaire}

\textbf{Anonymisation}~: Processus irréversible consistant à rendre impossible toute identification d'une personne à partir de données personnelles.

\textbf{API REST}~: Interface de programmation applicative respectant les principes architecturaux REST (Representational State Transfer), utilisant les méthodes HTTP (GET, POST, PUT, DELETE).

\textbf{Bcrypt}~: Algorithme de hachage de mots de passe adaptatif, conçu pour être intentionnellement lent afin de résister aux attaques par force brute.

\textbf{CLAHE}~: Contrast Limited Adaptive Histogram Equalization. Technique d'amélioration de contraste local évitant la surexposition des zones déjà contrastées.

\textbf{CLIP}~: Contrastive Language--Image Pre-training. Modèle de vision-langage d'OpenAI entraîné sur 400 millions de paires image-texte, capable de classification zéro-shot.

\textbf{DICOM}~: Digital Imaging and Communications in Medicine. Standard international pour le stockage, l'échange et la transmission d'images médicales numériques.

\textbf{Deep Learning}~: Sous-domaine du Machine Learning utilisant des réseaux de neurones à de nombreuses couches pour apprendre des représentations hiérarchiques des données.

\textbf{FastAPI}~: Framework Python web moderne, haute performance, pour la construction d'APIs avec Python 3.7+ basé sur les annotations de type standard de Python.

\textbf{Inpainting}~: Technique de reconstruction numérique consistant à remplir des régions manquantes ou masquées d'une image en propageant les valeurs des pixels environnants.

\textbf{IoU (Intersection over Union)}~: Métrique de chevauchement entre deux boîtes englobantes, calculée comme le rapport de l'intersection sur l'union des deux surfaces.

\textbf{JWT (JSON Web Token)}~: Standard RFC~7519 définissant un moyen compact et autonome de transmettre des informations entre parties sous forme de JSON signé numériquement.

\textbf{MinIO}~: Serveur de stockage objet open-source haute performance, compatible avec l'API Amazon S3, déployable sur site.

\textbf{OCR (Optical Character Recognition)}~: Reconnaissance optique de caractères. Technologie permettant de localiser et d'identifier le texte présent dans une image numérique.

\textbf{PaddleOCR}~: Bibliothèque OCR open-source de Baidu, basée sur PaddlePaddle, offrant des performances de pointe pour la détection et la reconnaissance de texte multilingue.

\textbf{PHI (Protected Health Information)}~: Selon HIPAA, toute information de santé individuellement identifiable relative à l'état de santé passé, présent ou futur d'un individu.

\textbf{RBAC (Role-Based Access Control)}~: Contrôle d'accès basé sur les rôles. Modèle de sécurité dans lequel les autorisations sont associées à des rôles et les utilisateurs se voient attribuer des rôles.

\textbf{RGPD}~: Règlement Général sur la Protection des Données. Réglementation européenne (2016/679) sur la protection des données à caractère personnel, entrée en vigueur le 25 mai 2018.

\textbf{TELEA}~: Algorithme d'inpainting basé sur la méthode Fast Marching, proposé par Alexandru Telea en 2004 et implémenté dans OpenCV sous le nom \texttt{cv2.INPAINT\_TELEA}.

\textbf{ViT (Vision Transformer)}~: Architecture de réseau de neurones appliquant le mécanisme d'attention multi-têtes (Transformer) directement sur des séquences de patches d'images.

\textbf{Zéro-shot}~: Capacité d'un modèle à classifier ou reconnaître des catégories qu'il n'a jamais vues explicitement pendant l'entraînement, en s'appuyant sur des descriptions textuelles.
